This chapter formulates the thesis problem from three complementary perspectives: (i) task-space generalization under sparse demonstrations, (ii) kinematic motion generation as conditional sequence modeling, and (iii) constrained planning for long-horizon rollout. Each perspective captures a different aspect of the same objective: generating reusable whole-body motion plans for object relocation.

\section{Task Space and Sparse-Demonstration Setting}

We consider object-relocation tasks in which the desired box displacement is a vector
\[
g \in \mathcal{G} \subseteq \mathbb{R}^{3},
\]
typically representing local-frame translation $(\Delta x,\Delta y,\Delta z)$. In practice, only a finite set of demonstrated goals is available:
\[
\mathcal{G}_{\mathrm{demo}} = \{g^{(1)},\dots,g^{(N)}\} \subset \mathcal{G}, \qquad N \ll \infty.
\]
The key difficulty is that $\mathcal{G}$ is continuous and large, whereas $\mathcal{G}_{\mathrm{demo}}$ is sparse. Hence the planner must interpolate and extrapolate to unseen goals $g \notin \mathcal{G}_{\mathrm{demo}}$.

In the full task definition, object relocation is naturally a 6D objective (3D translation and 3D orientation change). In this thesis, the planning target is simplified to a 3D displacement space to isolate and study kinematic trajectory generation under sparse goal coverage while keeping the analysis tractable.

An important representation decision is yaw invariance: trajectories that are equivalent up to global heading should map to nearby feature trajectories. Let $\psi_0$ be the anchor yaw at window start, and let $R_z(-\psi_0)$ be the planar rotation into the anchor frame. For any world-frame planar displacement $\Delta p_t^{w}\in\mathbb{R}^2$, the yaw-normalized displacement is
\[
\Delta p_t^{loc}=R_z(-\psi_0)\,\Delta p_t^{w}.
\]
This reduces effective sparsity by collapsing orientation-equivalent demonstrations into a shared local coordinate system.

\begin{figure}[!h]
    \centering
    \includegraphics[width=0.78\linewidth]{yaw_invariance.png}
    \caption{Yaw-invariant trajectory encoding. Global orientation differences are normalized in the local anchor frame, improving data reuse across sparse demonstrations.}
    \label{fig:problem_yaw_invariance}
\end{figure}

Let $s_t \in \mathcal{S} \subset \mathbb{R}^{d}$ denote the robot--object kinematic state at time $t$, including base pose terms, body-joint coordinates, and object-relative channels. Given a history window
\[
h_t = (s_{t-H_h+1},\dots,s_t),
\]
and a task goal $g$, the planner outputs a future trajectory segment
\[
\tau_{t+1:t+H} = (s_{t+1},\dots,s_{t+H}).
\]

\section{Kinematic Motion Generation with Reusable Primitives}

Rather than memorizing full demonstrations end-to-end, we treat each demonstration as a composition of shorter reusable motion primitives. Intuitively, different combinations of lower-body locomotion, torso orientation, and upper-body manipulation phases are required to cover the full task space of object displacements.

Let a full trajectory be represented as a sequence of primitive segments:
\[
\tau = \Pi(\pi_1,\pi_2,\dots,\pi_M),
\]
where each primitive $\pi_m \in \mathbb{R}^{H\times d}$ is a short kinematic segment and $\Pi$ denotes temporal composition. Generalization is then achieved by recombining locally valid segments under new task conditions, instead of reproducing one fixed global template.

This view motivates segmented training windows and autoregressive composition at inference: the model learns a local primitive distribution and reuses it to synthesize unseen long-horizon trajectories.

\begin{figure}[t]
    \centering
    \includegraphics[width=0.72\linewidth]{trajectory_segmentation.png}
    \caption{Kinematic trajectory segmentation into short reusable windows. The planner learns local motion primitives and composes them for long-horizon generation.}
    \label{fig:problem_segmentation_primitives}
\end{figure}

\section{Learning Objective as Conditional Distribution Modeling}

Because multiple physically plausible strategies can solve the same relocation task, the target is inherently multi-modal. We therefore model a conditional trajectory distribution
\[
p_{\theta}(\tau \mid h_t,g)
\]
instead of a deterministic map.

From demonstrations, we obtain a finite dataset
\[
\mathcal{D}=\{(h_i,g_i,\tau_i)\}_{i=1}^{M}, \qquad g_i \in \mathcal{G}_{\mathrm{demo}}.
\]
In our setting, these expert trajectories are not directly captured from robot execution logs. Instead, they come from a data-generation pipeline that first retargets human demonstrations into kinematically feasible humanoid motions and then refines them with sampling-based trajectory optimization to improve dynamic consistency. This design follows the progression from interaction-preserving retargeting in OmniRetarget to trajectory-level dynamic refinement in DynaRetarget \cite{omniretarget,dynaretarget}.

The statistical goal is to minimize expected loss under the (unknown) data distribution:
\[
\theta^{\star} = \arg\min_{\theta}\; \mathbb{E}_{(h,g,\tau)\sim p_{\mathrm{data}}}\big[\ell(\tau, p_{\theta}(\cdot\mid h,g))\big],
\]
which is implemented through empirical risk minimization on $\mathcal{D}$.

Equivalently, one may view learning as conditional density matching:
\[
\min_{\theta}\; D_{\mathrm{KL}}\!\left(p_{\mathrm{data}}(\tau\mid h,g)\;\|\;p_{\theta}(\tau\mid h,g)\right).
\]
Generalization quality is determined by how well $p_{\theta}(\tau\mid h,g)$ extends to sparse and unseen goals.

The formulation also supports iterative data augmentation. Let $\mathcal{T}$ denote a tracking operator (RL controller plus environment rollout) and $\mathcal{F}$ a validity filter that accepts trajectories satisfying task and kinematic-quality criteria. Starting from demonstrations $\mathcal{D}^{(0)}$, one iteration is
\[
\tilde{\mathcal{D}}^{(r)}=\mathcal{F}\!\left(\mathcal{T}(\tau),\; \tau\sim p_{\theta^{(r)}}(\tau\mid h,g)\right),
\qquad
\mathcal{D}^{(r+1)}=\mathcal{D}^{(r)}\cup\tilde{\mathcal{D}}^{(r)}.
\]
This defines an evolutionary loop in which generated-and-tracked motions expand training coverage over the goal space, and the diffusion planner is re-estimated on the augmented dataset, following the same high-level principle as iterative generator--tracker pipelines \cite{xu2025parc}.

\section{Constrained Planning View}

For a sampled trajectory $\tau=(s_{1},\dots,s_{H})$, feasible behavior should satisfy three conditions.

\paragraph{Task completion.}
Let $o_t=\phi(s_t)$ denote the object pose component. For a desired displacement $g$, endpoint error should remain below a tolerance:
\[
\|\psi(o_H,o_0)-g\|_2 \le \varepsilon_g,
\]
where $\psi(o_H,o_0)$ extracts achieved displacement.

\paragraph{Kinematic admissibility.}
State and incremental motion should remain in admissible sets:
\[
s_t \in \mathcal{S}_{\mathrm{adm}}, \qquad (s_{t+1}-s_t)\in \mathcal{V}_{\mathrm{adm}}.
\]
This compactly captures joint limits, plausible frame-to-frame variation, and representation-level sanity constraints.

\paragraph{Demonstration-manifold consistency.}
Generated trajectories should remain close to the demonstrated behavior manifold, avoiding implausible yet numerically valid states.

\paragraph{Conceptual objective.}
A compact objective view is
\[
\min_{\tau}\; J_{\mathrm{task}}(\tau;g)+\lambda_{\mathrm{kin}}J_{\mathrm{kin}}(\tau)+\lambda_{\mathrm{prior}}J_{\mathrm{prior}}(\tau),
\]
where diffusion learning provides the motion prior implicit in $J_{\mathrm{prior}}$.

At the pipeline level, the objective extends across iterations as
\[
\min_{\theta,\,\pi}\;\mathbb{E}_{g\sim p(g)}\!\left[J\!\left(\mathcal{T}_{\pi}(\tau),g\right)\right],
\qquad \tau\sim p_{\theta}(\tau\mid h,g),
\]
with planner parameters $\theta$ and tracker policy $\pi$ improved through alternating planner training and tracker-based data augmentation.

\section{Conditional Diffusion for Long-Horizon Trajectory Planning}
In this thesis, conditional diffusion is used as a trajectory-window generator inside a long-horizon planner. The model predicts future kinematic segments under state-history and goal conditions, while rollout over long tasks is obtained through autoregressive stitching of successive windows. The detailed denoising equations, $v$-prediction objective, and stitching mechanism are deferred to Chapter 4, where they are presented as methodological design choices.

\subsection{Evaluation axes implied by the formulation}

The above formulation implies three complementary evaluation axes:
\begin{itemize}
    \item \textbf{In-distribution fidelity:} Does $p_{\theta}(\tau\mid h,g)$ reproduce demonstrated local motion patterns?
    \item \textbf{Task-space generalization:} Does performance remain strong for goals in $\mathcal{G}\setminus\mathcal{G}_{\mathrm{demo}}$?
    \item \textbf{Constraint consistency:} Do sampled trajectories remain kinematically plausible under long-horizon stitching?
\end{itemize}

These criteria define the empirical analyses in later chapters and directly measure whether sparse-demonstration learning can scale to a large continuous displacement space in $\mathbb{R}^3$.
